% \iffalse meta-comment
%
%% File: latex-lab-graphic.dtx (C) Copyright 2022-2025 LaTeX Project
%
% It may be distributed and/or modified under the conditions of the
% LaTeX Project Public License (LPPL), either version 1.3c of this
% license or (at your option) any later version.  The latest version
% of this license is in the file
%
%    https://www.latex-project.org/lppl.txt
%
%
% The development version of the bundle can be found below
%
%    https://github.com/latex3/latex2e/required/latex-lab
%
% for those people who are interested or want to report an issue.
\def\ltlabgraphicdate{2025-07-27}
\def\ltlabgraphicversion{0.80h}
%
%<*driver>
\DocumentMetadata{tagging=on,pdfstandard=ua-2}
\documentclass{l3doc}
\usepackage{latex-lab-testphase-l3doc}
\EnableCrossrefs
\CodelineIndex
\begin{document}
  \DocInput{latex-lab-graphic.dtx}
\end{document}
%</driver>
%
% \fi
%
%
% \title{The \textsf{latex-lab-graphic} package\\
% Tagging of included graphics }
% \author{\LaTeX{} Project\thanks{Initial implementation done by Ulrike Fischer}}
% \date{v\ltlabgraphicversion\ \ltlabgraphicdate}
%
% \maketitle
%
% \newcommand{\xt}[1]{\textsl{\textsf{#1}}}
% \newcommand{\TODO}[1]{\textbf{[TODO:} #1\textbf{]}}
% \newcommand{\docclass}{document class \marginpar{\raggedright document class
% customizations}}
%
% \providecommand\hook[1]{\texttt{#1}}
% \NewDocElement[printtype=\textit{plug},idxtype=plug,idxgroup=Plugs]{Plug}{plugdecl}
% \NewDocElement[printtype=\textit{socket},idxtype=socket,idxgroup=Sockets]{Socket}{socketdecl}
%
% \newcommand\tikzname{ti\textit{k}Z}
% \begin{documentation}
% \begin{abstract}
% The following code implements a first draft for the tagging of graphics
% included with \cs{includegraphics} and the \env{picture} environment.
% \end{abstract}
% 
% 
% \section{Introduction}
% 
% Tagging of pictures is non trivial.
%
% \begin{enumerate}
% \item Pictures generally can have various purposes and so different \emph{tagging modes} must
% be provided:
% \begin{itemize}
% \item The pictures can be purely \emph{ornamental} and \emph{decorative}, 
% e.g., (always/usually) some page border.
% This should normally be tagged as \emph{artifact}, either in a \texttt{artifact} 
% MC-chunk or as an \texttt{Artifact} structure\footnote{The second option exists only in PDF 2.0}.
%
% \item The pictures can be \emph{illustrative figures}.
% This should normally be tagged as a \texttt{Figure} structure with \emph{alternative text}.
%
% \item The pictures can represent a symbol. This should be tagged as a \texttt{Span} structure element
% with an \texttt{/ActualText} mapping to some Unicode Codepoint(s) 
% (or perhaps even directly in the stream with a Span BDC with an \texttt{/ActualText}).
%
% \item The pictures can also be intended for consumption as normal \emph{text}
% For example the todonotes package uses a
% \tikzname{} picture to surround the text in a node with some colored frame and 
% background.
% In this case the text (which can contain more elements like lists)
% should be tagged as usual and put in an \texttt{Aside} structure while the decorative
% elements should be marked up as artifacts. 
% 
% \item There are also case where no tagging is wanted, e.g. because the tagging
% is done by surrounding commands, in this case the begin/end sockets should be transparent and
% do nothing.
%  
% \item And naturally there can be more complicated scenarios with
% mixtures of these elements, e.g., some text 
% in nodes in a \tikzname{} picture can be meaningful 
% while other nodes still should be tagged as artifact.
% \end{itemize}
%
% \item The various packages that allow to draw pictures uses lots of boxes and 
% moves them around and that makes is not easy to get the tagging right --
% especially with pdflatex where one has to insert the literals at the right time.
%
% \item When the picture is tagged as \texttt{Figure}, the PDF-UA standards require
% an attribute with a \texttt{BBox} key in the structure. That BBox describes 
% the placement of the picture on the page, 
% and this typically require some low-level hacking into the picture code to calculate
% the values.
% \end{enumerate}
% 
% The code here handles directly only the tagging of pictures included with \cs{includegraphics}
% and the \env{picture} environment. But the used method and the documentation tries
% to stay as general as possible, to make it easy to adapt the code to pictures drawn
% with other packages like \pkg{l3draw}, \pkg{tikz}, \pkg{pstricks}, \pkg{luamplibs} or similar packages. 
% 
% \section{General implementation needs and ideas}
% 
% \subsection{User interfaces}
% 
% The tagging of pictures can not be done fully automatically: user
% have to add alternative text or mark up a picture as an artifact. 
% It is important that the relevant user interfaces are similar across 
% all picture/graphic packages to make it
% easy for authors to adapt their documents. 
% 
% We therefore recommend package authors to implement an interface 
% a key-value interface like the following. 
% 
% \subsubsection{Tagging mode of individual graphics}
% To change the tagging mode of an individual graphic, the keys
% \texttt{artifact} (decorations), \texttt{alt} (illustrative), 
% and \texttt{actualtext} (symbol) should be used. The keys \texttt{alt}
% (and also \texttt{actualtext}) take a text as argument. 
% 
% \begin{verbatim}
% \includegraphics[artifact]{example-image}
% \begin{tikzpicture}[artifact] ... \end{tikzpicture}
% \tikz[artifact]{...}
% \begin{picture}[artifact]...\end{picture}
% 
% \includegraphics[alt=example image]{example-image}
% \begin{tikzpicture}[alt=\myalttext] ... \end{tikzpicture}
% \tikz[alt=example image]{...}
% \begin{picture}[alt=example image]...\end{picture}
% 
% \includegraphics[actualtext=A]{example-image-A}
% \begin{tikzpicture}[actualtext=A] ... \end{tikzpicture}
% \tikz[actualtext=A]{...}
% \begin{picture}[actualtext=A]...\end{picture}
% \end{verbatim}
% 
% Additionally, a key \texttt{tagging-setup} can be 
% provided that allows to handle more complicated cases, e.g.,
%  
% \begin{verbatim}
% \includegraphics[tagging-setup={false}]{example-image} %no tagging at all
% \begin{tikzpicture}[tagging-setup={artifact}] ... \end{tikzpicture}
% \begin{tikzpicture}[tagging-setup={text}] ... \end{tikzpicture}
% \begin{tikzpicture}[tagging-setup={alt=Water (H20),tag=Formula}] ... \end{tikzpicture}
% \end{verbatim}
% 
% This key is fully implemented for \cs{includegraphics} and \env{picture} but only
% partially in the \tikzname{} module.
% 
% \subsubsection{Setting the tagging mode for a scope}\label{sec:scope}
% 
% Packages can also provide interfaces to change the tagging mode for all pictures or
% nodes in a scope. 
% 
% If this is done side-effects on the tagging of other picture 
% types must be considered: 
% If the generic sockets declared below are used 
% (which are then used also by \cs{includegraphics}, \env{picture}
% and perhaps more environments) a change of the tagging mode can
% affect all pictures types using these sockets in the same scope. 
% 
% This means that packages that want to offer \enquote{scope support} 
% but restrict it to their own picture type should either use 
% their own sockets modelled after the generic sockets, 
% or use some specific variable to control the change of the tagging mode. 
% 
% With \tikzname{} \enquote{scope support} could be implemented rather easily as it has a setup command anyway.
% So with the implementation in latex-lab-tikz, all these work as expected
% \begin{verbatim}
% \tikzset{artifact}
% \tikzset{alt=a text}
% \tikzset{actualtext=A}
% \tikzset{tagging-setup=text}
% \end{verbatim}
%  
% With \cs{includegraphics} the standard \cs{setkeys} can be used:
% \begin{verbatim}
% \setkeys{Gin}{alt=a text}
% \setkeys{Gin}{artifact}
% \setkeys{Gin}{actualtext=A}
% \end{verbatim}
% This will then also affect \env{picture} environments in the scope.
% 
% \subsection{Default tagging mode}
% 
% If none of the keys are used, the picture code must chose a default tagging mode.  
% The best one depends on the graphic type and 
% should be chosen as needed and then documented.
% 
% With \cs{includegraphics} we use as default the illustrative mode,
% tag it as \texttt{Figure} structure, use the file name as alternative text
% and issue a warning that a real alternative text is missing.
% 
% With the \env{picture} enviroment we use as default the illustrative mode
% to, but use a fix text as alternative text and issue a warning too.
% 
% With \tikzname{} (a first implementation is currently in latex-lab-tikz), 
% the default is the text mode, which means that a screen reader will read
% the text in the nodes in the order they appear in the code. With this default,
% e.g. \cs{todo}'s from the \env{todonotes} are correctly tagged.  
% 
%  
% In this implementation the \texttt{artifact}
% key can also be applied on nodes and will remove them from the tagging.
% \begin{verbatim}
% \begin{tikzpicture}
% \node[draw=red](x){Important!};
% \node[artifact,fill=blue,anchor=west] at (x.east){\phantom{Ip}};
% \end{tikzpicture}
% \end{verbatim}

% \subsection{Sockets, plugs and commands}
% 
% \begin{description}
% \item[sockets]
%  The example implementations make use of up-to five \emph{tagging sockets}.
%  The sockets are all used by the various graphic codes inside a group:
%  \begin{itemize}
%  \item  An initialization (\texttt{init}) socket that handles 
%  the key-val argument and setups the tagging mode by switching 
%  the plugs of the other tagging sockets. 
%  \item  Two sockets that are used at the \texttt{begin} and \texttt{end} of 
%  the picture and setup the main structure element. 
%  The \texttt{end} socket typically also executes if needed the code to calculate the BBox
%  attribute and add it to the structure.
%  After the \texttt{begin} socket tagging is suspended, and before the
%  \texttt{end} socket resumed.
%  \item Two sockets for the text mode that are used around places where a 
%  picture contains text. Tagging must be resumed before this sockets
%  if they should do anything at all.
%  \end{itemize}
%  
%  The sockets take arguments that allows some configuration (e.g. to use
%  a special command for the BBox calculation) and the plugs are coded so that
%  they can be used in more than one picture type (they are shared here
%  between \cs{includegraphics} and \env{picture}) but there is no obligation 
%  to use them in every graphic code. 
%  A package can declare its own sockets and plugs. See \ref{sec:scope} for some discussion 
%  why this can be a good idea.
%  
% \item[plugs] 
% The main \texttt{begin} and \texttt{end} sockets should normally have 
% four \emph{plugs} 
% for the different modes: 
% \texttt{alt}, \texttt{actualtext}, \texttt{artifact} and \texttt{text}.
% To disable tagging, \cs{SuspendTagging} can be used,
% alternatively for sockets with zero or one argument, 
% the predefined \texttt{noop} plug can be assigned,
% and for sockets with two arguments (if the socket has been declared with
% \cs{NewTaggingSocket} the \texttt{transparent} 
% socket which lets pass through the 
% second argument.
%  
% The initialization socket and the texts sockets typically have only one 
% additional \emph{plug} that is used when tagging is active. 
% 
% \item[commands] A command to store the position of a reference 
% point on the page and a command to calculate the BBox from this reference
% point and the size of the picture are needed. The second command must also add the 
% attribute to the main structure element (how to do this can be seen in the implementation).
% These commands are typically specific for a picture type.
% 
%\end{description}
% 
% We now describe how this general principles have be implemented 
% in the concrete examples of the \cs{includegraphics} command, the \env{picture} environment, 
% TODO! and a simple environment for l3draw commands.
% 
% \section{Sockets, plugs, commands}
% 
%  The code defines the following generic sockets and commands.
%  
%  \begin{itemize} 
%  \item \texttt{graphic/init} with the plug \texttt{default}
%  \item \texttt{graphic/begin} (one argument) with the plugs \texttt{alt},
%   \texttt{actualtext}, \texttt{artifact}, \texttt{text}, \texttt{off}. The argument
%   allows to set a default alternative text.
%  \item \texttt{graphic/end} (two arguments) with the plugs \texttt{alt},
%   \texttt{actualtext}, \texttt{artifact}, \texttt{text}, \texttt{off}. The first argument
%   of the socket typically receives the command to calculate the BBox. This command
%   often contains a savepos command and so has to go before the graphic box. The second
%   argument allows to insert the box between this command and the tagging commands.%   
%  \item \texttt{graphic/text/begin} no argument, with the plug \texttt{default} 
%  \item \texttt{graphic/text/end} no argument with the plug \texttt{default}
%  \end{itemize}
%  
%  \begin{function}{\l_tag_graphic_mode_tl}
%  This variable holds the current active graphic mode. It is e.g. used
%  to test if text should resume tagging.
%  \end{function}
% 
% \section{Make \cs{includegraphics} tagging aware}
% 
% Tagging of graphics included with \cs{includegraphics} is at a first glance
% easier to handle than, e.g., a \tikzname{} graphic, as there is only a simple
% box with a picture and no text content to consider. 
% One would think that adding some structure commands around a box shouldn't pose much problems. 
%  
% But things are actually not so easy. 
% 
% At first such graphics are inserted into the PDF as XObjects and
% there are two ways to add an such an XObject to a structure: 
% similar to text as a marked content item (by surrounding 
% it with \cs{tagmcbegin} and \cs{tagmcend}) or by referencing the XObject 
% with an OBJR object (similar to a link annotation). Which method is more sensible 
% (and if it actually matters) is unknown. 
% Currently the first method is used as the second would changes in the backend files.
%   
% At second---and this is actually a \emph{much} bigger problem\footnote{Which shows
% also in the amount of code dedicated here only to this problem.}---if the graphic
% is tagged as an illustrative picture the 
% \texttt{Figure} structure element should have an attribute with an \texttt{BBox} entry.
% The value of this \texttt{BBox} is an array of four numbers that gives the 
% coordinates of the left, bottom, right, and top edges 
% of the structure element’s bounding box on the page. 
% That is the rectangle that completely encloses
% its \emph{visible} content and so has not necessarily the same size as the TeX bounding box: 
% if \texttt{viewport} or \texttt{trim} is used and the
% graphic is not clipped, the visible content can be larger. 
% It turned out to be extremely tricky to get a sensible result, and there are still
% open problems and restrictions.
% 
%  
% \subsection{The BBox calculation}
% The reference point on the page is retrieved with \verb+\tex_savepos:D+ and a property
% just before the graphic.
%  
% Getting from this the \texttt{BBox} can be quite straightforward for a graphic that is 
% used once as is. But graphics can be trimmed, scaled, reflected, rotated and reused
% in various ways. These transformations typically involve a mix of \TeX{} commands
% that shift a box or change the bounding box and backend commands that insert
% a pdfliteral with a transformation matrix. Calculating
% the correct \texttt{BBox} in all cases is not possible 
% without rewriting large parts of the graphics and graphicx packages. Problematic are
%  \begin{itemize}
%  \item manipulations through external box commands 
%  (\cs{rotatebox}, \cs{reflectbox}, \cs{scalebox}). Their implementation
%  do not pass the transformation matrix in a way 
%  that allows to track the changes for the \texttt{BBox} 
%  of an included graphic: sometimes
%  the values are set to late (after the box is already stored), and often the
%  values are not grouped and can leak out from earlier uses of the commands.
%  
%  \item some combination of keys in the optional argument of \cs{includegraphics}.
%  Examples are \texttt{origin} and multiple calls to \texttt{scale} 
%  and \texttt{angle}) as they internally call the box commands.
%  Examples of failing combinations can be found in the test file 
%  \texttt{graphic-faults}.
%  
%  \item graphics that are stored in a box and reused: to get the \texttt{BBox} 
%  one has to set a label that stores the position with \cs{pdfsavepos}, 
%  and if a box is reused one gets multiply defined labels. 
%  One possible solution here is to make
%  use of the new delayed \cs{pdfliteral}. It allows to change the label names
%  in the shipout, but this requires careful tracking the box usages 
%  and so various kernel changes.
%  \end{itemize}
%  
%  
%  Therefore a correct BBox is currently 
%  implemented only for simple \cs{includegraphics} and the keys
%  \texttt{viewport}, \texttt{trim}, 
%  \texttt{scale} and \texttt{angle} (used at most once). 
%  
%  Currently not supported are
%  \begin{itemize}
%  \item graphics inside \cs{rotatebox}, \cs{reflectbox}, \cs{scalebox}. 
%  
%  TODO: A new implementation with \texttt{l3graphics} and \texttt{l3box} is 
%  probably needed here.
%  
%  \item  multiple uses of the \texttt{scale} and \texttt{angle} keys
%   
%  \item multiple use of graphics stored in boxes. For such graphics automated tagging
%  should be probably deactivated when storing the content 
%  and tagging should be added around the \cs{usebox}. (How to 
%  proceed when content is saved in boxes needs generally more testing).    
%  \end{itemize}
%  
%  
%  \subsection{User interface}
%  As suggested above the code (re)defines keys for \cs{includegraphics}
%  to add the recommanded interfaces \texttt{alt}, \texttt{artifact} and \texttt{actualtext}.
%  It also offers some keys specific to \cs{includegraphics}:
%  \begin{description}
%  \item[\texttt{alt}] This key is already defined in the graphicx package
%  but redefined here to switch to the illustrative mode and to add its value
%  as alternative text. 
%  \item[\texttt{actualtext}] This switches to the actualtext mode.
%  This is useful for small graphics that represent single chars or a short word like a 
%  logo. If \texttt{actualtext} is used, the graphics is not enclosed in \texttt{Figure}
%  structure but in a \texttt{Span} structure and no \texttt{/BBox} attribute is added.%
%  \footnote{This in accordance with (the draft of) PDF/UA-2 but violates perhaps PDF/UA-1.}
%  \item[\texttt{artifact}] This tags the graphic as an artifact.
%  \item[\texttt{tagging-setup}] This key takes as argument a key-list,
%  which can contain the following keys:
%   \begin{description}
%    \item[\texttt{alt=\meta{text}}] This a second way to tag the graphic as figure with
%    alternative text. \texttt{alt} without value will tag as figure but not change
%    the text variable, \texttt{alt=} will empty the text variable (and typically trigger
%    a warning).
%    \item[\texttt{actualtext=\meta{text}}] This a second way to tag the graphic as a
%    symbol with actualtext.  
%    \item[\texttt{artifact}] This a second way to tag the graphic as artifact. 
%    \item[\texttt{text}] This switches to text mode.
%    \item[\texttt{off}] When used tagging will be stopped completely. It is then the 
%   responsibility of the surrounding code to add appropriate tagging commands.
%   \item[\texttt{tag=\meta{name}}] This switches to the illustrative mode but uses
%   \meta{name} as tag name in the structure instead of the default \texttt{Figure}. 
%   This can for example be used to tag an image of a formula with \texttt{Formula}.
%   \end{description}   
%  \item[\texttt{adjust-BBox}] If the calculated \texttt{/BBox} values are wrong
%  they can be corrected with this key. It expects four dimensions that are added to the 
%  \texttt{/BBox} values. 
%  \end{description}
%  
%  The code also add to the \texttt{debug} key of \cs{DocumentMetadata} 
%  the value \texttt{BBox}.
%  This adds a half transparent red layer showing the calculated BBox.
%  \begin{verbatim}
%   \DocumentMetadata{tagging=on,debug=BBox}
%  \end{verbatim}
%  
%  \section{Make \env{picture} tagging aware}
%  
%  \subsection{User interface}
%  The \env{picture} environment doesn't have a optional argument with key-value value processing.
%  So the code changes that and then provides the keys already discussed for \cs{includegraphics}.
%  
%  \subsection{BBox calculation}
%  
%  The BBox calculation is much simpler than for \cs{includegraphics} as the code simply takes
%  the declared size from the \env{picture} arguments. 
%  
%  \subsection{Tagging in text mode}
%  
%  A  \env{picture} can contain \cs{put} commands with text content, tagging this in 
%  text mode could make sense in some cases. But it is not quite clear if one can/should 
%  redefine the \cs{put} command. If text tagging is wanted we suggest to 
%  define a dedicated command along these lines:
%  \begin{verbatim}
%  \NewDocumentCommand\picturenode{O{}r()m}
%    {
%     \group_begin:
%     \keys_set:nn{tag/graphic}{#1}
%     \str_if_eq:VnT\l_tag_graphic_mode_tl {text}
%      {\tag_resume:n{\picturenode}}
%     \tag_socket_use:n{graphic/text/begin}
%     \put(#2){#3}
%     \tag_socket_use:n{graphic/text/end}  
%     \group_end:
%    }   
%  
%  \end{verbatim}   
%  
%  \section{Make an l3draw environment tagging aware}
%  The newest l3draw version (in the latex3 github)
%  has all needed data to define a command to retrieve the BBox
%  and to build a tagging aware environment. l3draw currently has
%  no dedicated function to add text, instead one has to store the text
%  in a box and then reuse it, so similar to the \env{picture} environment
%  a dedicated command with tagging awareness is suggested.
%  
%  A command to calculate the BBox can for example be defined like this
%  \begin{verbatim}
%  \cs_new:Npn\draw_tag_bbox_attribute:
%   {
%     \tl_set:Ne \g__tag_graphic_lx_tl
%       {
%         \dim_to_decimal_in_bp:n 
%           { 
%             \property_ref:een {draw.\int_use:N\g_draw_id_int}{xpos}{0}sp 
%             + 
%             \g_draw_bb_xmin_dim 
%           }
%       }   
%     \tl_set:Ne \g__tag_graphic_ly_tl
%       {
%         \dim_to_decimal_in_bp:n 
%           { 
%             \property_ref:een {draw.\int_use:N\g_draw_id_int}{ypos}{0}sp 
%             + 
%             \g_draw_bb_ymin_dim 
%           }
%       }
%     \tl_set:Ne \g__tag_graphic_ux_tl
%       {
%         \dim_to_decimal_in_bp:n 
%           { 
%             \property_ref:een {draw.\int_use:N\g_draw_id_int}{xpos}{0}sp 
%             + 
%             \g_draw_bb_xmax_dim 
%           }
%       }   
%     \tl_set:Ne \g__tag_graphic_uy_tl
%       {
%         \dim_to_decimal_in_bp:n 
%           { 
%             \property_ref:een {draw.\int_use:N\g_draw_id_int}{ypos}{0}sp 
%             + 
%             \g_draw_bb_ymax_dim 
%           }
%       }          
%     \bool_if:NT\l__tag_graphic_debug_bool
%      {
%        \__tag_graphic_show_bbox:VVVVne
%         \g__tag_graphic_lx_tl
%         \g__tag_graphic_ly_tl
%         \g__tag_graphic_ux_tl
%         \g__tag_graphic_uy_tl
%         {red}
%         {draw.\int_use:N\g_draw_id_int}              
%      }
%    \tag_struct_gput:ene
%         {\tag_get:n{struct_num}}
%         {attribute}
%         {
%           /O /Layout /BBox~
%           [
%             \dim_to_decimal_in_bp:n 
%              { 
%               \property_ref:een {draw.\int_use:N\g_draw_id_int}{xpos}{0}sp 
%                + 
%               \g_draw_bb_xmin_dim 
%              }
%             \c_space_tl 
%             \dim_to_decimal_in_bp:n 
%              {              
%                \property_ref:een {draw.\int_use:N\g_draw_id_int}{ypos}{0}sp  
%                +
%                \g_draw_bb_ymin_dim
%              }
%             \c_space_tl 
%             \dim_to_decimal_in_bp:n 
%              { 
%                \property_ref:een {draw.\int_use:N\g_draw_id_int}{xpos}{0}sp  
%                +
%                \g_draw_bb_xmax_dim
%              }
%             \c_space_tl 
%             \dim_to_decimal_in_bp:n 
%              { 
%                \property_ref:een {draw.\int_use:N\g_draw_id_int}{ypos}{0}sp  
%                +
%                \g_draw_bb_ymax_dim
%              }   
%           ]
%         }
%    }
% \end{verbatim}
% 
% A tagging aware environment can then be defined like this
% \begin{verbatim}
% \NewDocumentEnvironment{tagged-draw}{O{}}
%  {\leavevmode
%   \ExplSyntaxOn
%   \tag_socket_use:nn{graphic/init}{#1} 
%   \tag_socket_use:nn{graphic/begin}{tagged-draw~environment}
%   \tag_suspend:n{\draw} %
%   \draw_begin:\ignorespaces
%  }
%  {
%   \draw_end:
%   \tag_resume:n{\draw}
%   \tag_socket_use:nnn{graphic/end}{\draw_tag_bbox_attribute:}{} 
%  }
% \end{verbatim}
%  
% And a command for tagged text nodes could look like this
% \begin{verbatim}
% \cs_new_protected:Npn \draw_text_node:nnn #1#2#3 %#1 keyval, #2 point, #3 text
%  {
%    \group_begin:
%    \keys_set:nn{tag/graphic}{#1}
%    \str_if_eq:VnT\l_tag_graphic_mode_tl {text}      
%       {\tag_resume:n{\draw_node:nn}}
%    \tag_socket_use:n{graphic/text/begin}
%    \hbox_set:Nn \l_tmpa_box{#3}
%    \draw_box_use:Nn\l_tmpa_box {#2}
%    \tag_socket_use:n{graphic/text/end}
%    \group_end:    
%  }
% \end{verbatim}   
%    \begin{macrocode}
%<@@=tag>
%<*package>
%    \end{macrocode}
% \end{documentation}
% \begin{implementation}
% \section{Implementation} 
%    \begin{macrocode}
\ProvidesExplPackage {latex-lab-testphase-graphic} {\ltlabgraphicdate} {\ltlabgraphicversion}
  {Code related to the tagging of graphics}
%    \end{macrocode}
% a variant
%    \begin{macrocode}
\cs_generate_variant:Nn \tag_socket_use:nn {ne}
%    \end{macrocode}
% We load l3opacity for the debug code if opacity is not already defined in the kernel:
%    \begin{macrocode}
\cs_if_exist:NF \opacity_select:n
  {
    \RequirePackage{l3opacity}
  }  
%    \end{macrocode}
%
% \begin{macro}{\@@_graphic_savepos:n}
% this is the command which stores the position. Similar to 
% zref-savepos it uses two savepos commands 
% for the case that bidi changes the processing order.
%    \begin{macrocode}
\cs_new_protected:Npn\@@_graphic_savepos:n #1
 {
  \tex_savepos:D 
  \property_record:nn{#1}{xpos,ypos,abspage}
  \tex_savepos:D 
 }
\cs_generate_variant:Nn \@@_graphic_savepos:n {e}
%    \end{macrocode}
% \end{macro}
%
% \subsection{Variables}
% \begin{variable}
%  {
%   \l_@@_graphic_alt_tl,
%   \l_@@_graphic_actual_tl,
%   \l_@@_graphic_struct_tl,
%   \l_tag_graphic_mode_tl,
%  }
%  
% These variables are related to the tagging. 
% Variables for the alt text, the actualtext and the structure tag.
% The variable that holds the tagging mode is public,
% so that commands can test for it, and e.g. restart
% tagging in text mode. 
% 
%    \begin{macrocode}
\tl_new:N \l_@@_graphic_alt_tl
\tl_new:N \l_@@_graphic_actual_tl
\tl_new:N \l_@@_graphic_struct_tl
\tl_set:Nn\l_@@_graphic_struct_tl {Figure}
%    \end{macrocode}
% \end{variable}
% 
% \begin{variable}{\l_@@_graphic_debug_bool}
% A boolean for debug code
%    \begin{macrocode}
\bool_new:N \l_@@_graphic_debug_bool
%    \end{macrocode}
% \end{variable}
% 
% The rest of the variables are related to the BBox calculation in \cs{includegraphics}.
% \begin{variable}{\g_@@_graphic_int}
% This is used to get unique labels in the savepos code.
%    \begin{macrocode}
\int_new:N\g_@@_graphic_int
%    \end{macrocode}
% \end{variable}

% \begin{variable}
% {
%  \g_@@_graphic_lx_tl,
%  \g_@@_graphic_ly_tl,
%  \g_@@_graphic_ux_tl,
%  \g_@@_graphic_uy_tl,
%  \l_@@_graphic_bboxcorr_seq
%  \l_@@_graphic_bboxcorr_bool
% }
% This commands will hold the calculated BBox values. Local variables
% would probably work too, but global variables can be easier retrieved
% in tests and debugging code ... 
%    \begin{macrocode}
\tl_new:N \g_@@_graphic_lx_tl
\tl_new:N \g_@@_graphic_ly_tl
\tl_new:N \g_@@_graphic_ux_tl
\tl_new:N \g_@@_graphic_uy_tl
\seq_new:N\l_@@_graphic_bboxcorr_seq
\bool_new:N\l_@@_graphic_bboxcorr_bool
%    \end{macrocode}
% \end{variable}
% \begin{variable}{\l_@@_graphic_currentlabel_tl}
% This holds the label name of the savepos.
%    \begin{macrocode}
\tl_new:N \l_@@_graphic_currentlabel_tl
%    \end{macrocode}
% \end{variable}
%
%
% \begin{variable}
%  {
%    \l_@@_graphic_sin_fp 
%   ,\l_@@_graphic_cos_fp
%   ,\l_@@_graphic_scale_fp
%   ,\l_@@_graphic_lxly_fp
%   ,\l_@@_graphic_lxuy_fp
%   ,\l_@@_graphic_uxly_fp
%   ,\l_@@_graphic_uxuy_fp
%   ,\l_@@_graphic_ux_fp
%   ,\l_@@_graphic_ly_fp
%   ,\l_@@_graphic_lx_fp
%   ,\l_@@_graphic_uy_fp
%   ,\l_@@_graphic_trim_ux_fp
%   ,\l_@@_graphic_trim_ly_fp
%   ,\l_@@_graphic_trim_lx_fp
%   ,\l_@@_graphic_trim_uy_fp
%  }
% A bunch of fp-variables (we don't use tl-vars, 
% to avoid to have to take care about minus signs everywhere) 
%    \begin{macrocode}
\fp_new:N\l_@@_graphic_sin_fp 
\fp_new:N\l_@@_graphic_cos_fp
\fp_new:N\l_@@_graphic_lxly_fp
\fp_new:N\l_@@_graphic_lxuy_fp
\fp_new:N\l_@@_graphic_uxly_fp
\fp_new:N\l_@@_graphic_uxuy_fp
\fp_new:N\l_@@_graphic_ux_fp
\fp_new:N\l_@@_graphic_ly_fp
\fp_new:N\l_@@_graphic_lx_fp
\fp_new:N\l_@@_graphic_uy_fp
%    \end{macrocode}
% this holds the scale value. Either \cs{Gin@scalex} or (if that is !) \cs{Gin@scaley}
%    \begin{macrocode}
\fp_new:N\l_@@_graphic_scale_fp
%    \end{macrocode}
% the follow variables hold the four trim values (or the equivalent calculated 
% values if viewport is used.
%    \begin{macrocode}
\fp_new:N\l_@@_graphic_trim_ux_fp
\fp_new:N\l_@@_graphic_trim_ly_fp
\fp_new:N\l_@@_graphic_trim_lx_fp
\fp_new:N\l_@@_graphic_trim_uy_fp
%    \end{macrocode}
% \end{variable}
% 
% \subsection{Tagging sockets}
% The sockets can perhaps not be shared between \cs{includegraphics} and \env{picture}
% but for now we try to do it, with the exception of the init socket.
%
% The begin socket takes an argument to allow to pass some configuration.
% The end socket takes two argument to allow to calculate the BBox
% before outputting the box.
% \begin{socketdecl}{tagsupport/graphic/init,tagsupport/graphic/begin,
% tagsupport/graphic/end,tagsupport/graphic/text/begin,tagsupport/graphic/graphic/text/end}
%    \begin{macrocode}
\NewTaggingSocket{graphic/init}{1}
\NewTaggingSocket{graphic/begin}{1}
\NewTaggingSocket{graphic/end}{2}
\NewTaggingSocket{graphic/text/begin}{0}
\NewTaggingSocket{graphic/text/end}{0}
%    \end{macrocode}
%\end{socketdecl}
%
% \subsection{Tagging plugs}
% 
% \subsubsection{Initialization of the tagging mode}
% \begin{plugdecl}{default (tagsupport/graphic/init),default (tagsupport/graphic/init)}
%     \begin{macrocode}
\NewTaggingSocketPlug{graphic/init}{default}
  {
   \keys_set:nn{tag/graphic}{#1}
   \ExpandArgs{no}\AssignTaggingSocketPlug{graphic/begin}{\l_tag_graphic_mode_tl}  
   \ExpandArgs{no}\AssignTaggingSocketPlug{graphic/end}{\l_tag_graphic_mode_tl}  
  }
\AssignTaggingSocketPlug{graphic/init}{default}  
%    \end{macrocode}
% \end{plugdecl}
% 
% \subsubsection{Main begin and end sockets}
% \begin{plugdecl}{figure (tagsupport/graphic/begin),figure (tagsupport/graphic/end)}
% These plugs handle the graphic as a figure.
% Around the graphic is a \texttt{Figure} environment which will
% use an alt text given in the optional argument and internally tagging is suspended.
% The Bbox will be set (after the second compilation) to the size of the declared size
%    \begin{macrocode}
\msg_new:nnn { tag } { alt-text-missing }
 {
   Alternative~text~for~graphic~is~missing.\\
   Using~'#1'~instead.
 }
\NewTaggingSocketPlug{graphic/begin}{alt}
  {
    \tag_mc_end_push: 
    \tl_if_empty:NT\l_@@_graphic_alt_tl         
      {
        \msg_warning:nne{tag}{alt-text-missing}{#1}
        \tl_set:Ne\l_@@_graphic_alt_tl {\text_purify:n{#1}}
      }
    \tag_struct_begin:n
      {
        tag=\l_@@_graphic_struct_tl,
        alt=\l_@@_graphic_alt_tl,
      }
    \tag_mc_begin:n{}      
  }
\NewTaggingSocketPlug{graphic/end}{alt}
  {
     #1
     #2
     \tag_mc_end:
     \tag_struct_end:
     \tag_mc_begin_pop:n{}
  }  
%    \end{macrocode}
% \end{plugdecl}
%
% \begin{plugdecl}{actualtext (tagsupport/graphic/begin),actualtext (tagsupport/graphic/end)}
% This plug handles the picture as a symbol with an actualtext.
% It tags the content as a Span and expects an actualtext.
% Internally tagging is suspended.
%    \begin{macrocode}
\NewTaggingSocketPlug{graphic/begin}{actualtext}
  {
   \tag_mc_end_push:
   \tag_struct_begin:n{tag=Span,actualtext=\l_@@_graphic_actual_tl }
   \tag_mc_begin:n{}
  }
\NewTaggingSocketPlug{graphic/end}{actualtext}
  {
   #2
   \tag_mc_end:
   \tag_struct_end:
   \tag_mc_begin_pop:n{}
  }  
%    \end{macrocode}
% \end{plugdecl}
%
% \begin{plugdecl}{artifact (tagsupport/graphic/begin),artifact (tagsupport/graphic/end)}
% This plug handles the picture as an artifact, as decoration.
% So it is surrounded by an artifact MC and internal text does not restart tagging.
%    \begin{macrocode}
\NewTaggingSocketPlug{graphic/begin}{artifact}
  {
    \tag_mc_end_push:
    \tag_mc_begin:n{artifact}
  }
\NewTaggingSocketPlug{graphic/end}{artifact}
  {
    #2
    \tag_mc_end:
    \tag_mc_begin_pop:n{}
  }
%    \end{macrocode}
% \end{plugdecl}
%
% \begin{plugdecl}{text (tagsupport/graphic/begin),text (tagsupport/graphic/end)}
% This plug can be used for text tagging.
% It basically does the same as the artifact plug. The main reason that it
% exist is that it looks more consistend and
% that we can test for the plug name and so restart tagging
% in places where this is wanted. (tagging can not be resumed inside a tagging hook,
% so has to use some external method).
%    \begin{macrocode}
\NewTaggingSocketPlug{graphic/begin}{text}
  {
    \tag_mc_end_push:
    \tag_mc_begin:n{artifact}
  }
\NewTaggingSocketPlug{graphic/end}{text}
  {
    #2
    \tag_mc_end:
    \tag_mc_begin_pop:n{}
  }
%    \end{macrocode}
% \end{plugdecl}
% 
% \begin{plugdecl}{off (tagsupport/graphic/begin),off (tagsupport/graphic/end)}
% This plug can be used for text tagging.
% It basically does the same as the artifact plug. The main reason that it
% exist is that it looks more consistend and
% that we can test for the plug name and so restart tagging
% in places where this is wanted. (tagging can not be resumed inside a tagging hook,
% so has to use some external method).
%    \begin{macrocode}
\NewTaggingSocketPlug{graphic/begin}{off}{}
\NewTaggingSocketPlug{graphic/end}{off}{#2}
%    \end{macrocode}
% \end{plugdecl}

% 
%  By default we use the alt plugs
%    \begin{macrocode}
\AssignTaggingSocketPlug{graphic/begin}{alt}
\AssignTaggingSocketPlug{graphic/end}{alt}
%    \end{macrocode}
%
% \begin{plugdecl}{default (tagsupport/graphic/text/begin),
%  default (tagsupport/picture/text/end)}
% These sockets are used inside the text
% plugs and ends the previous mc and restarts it after the text.
% Not used by \cs{includegraphics}. Unclear if they can be used by
% the \env{picture} environment.
%    \begin{macrocode}
\NewTaggingSocketPlug{graphic/text/begin}{default}
  {
   \tag_mc_end:
   \tag_mc_begin:n{}
  }
\NewTaggingSocketPlug{graphic/text/end}{default}
  {
   \tag_mc_end:
   \tag_mc_begin:n{artifact}
  }
\AssignTaggingSocketPlug{graphic/text/begin}{default}
\AssignTaggingSocketPlug{graphic/text/end}{default}
%    \end{macrocode}
% \end{plugdecl}
% 
% \subsection{Generic keys}
% The keys are used by picture and \cs{includegraphics} (through the 
% \texttt{tagging-setup} key
%    \begin{macrocode}
\keys_define:nn{tag/graphic}
   {
     ,alt .code:n =  
        {                    
          \tl_if_empty:nF{#1}
           { \tl_set:Ne\l_@@_graphic_alt_tl{\text_purify:n{#1}}}
          \tl_set:Nn\l_tag_graphic_mode_tl{alt}
        }
     ,artifact .code:n = 
        {
          \tl_set:Nn\l_tag_graphic_mode_tl{artifact}
        }
     ,actualtext .code:n = 
        {
          \tl_if_empty:nF{#1}
           { \tl_set:Ne\l_@@_graphic_actual_tl{\text_purify:n{#1}} }          
          \tl_set:Nn\l_tag_graphic_mode_tl{actualtext}
        }
     ,text .code:n =
        {
          \tl_set:Nn\l_tag_graphic_mode_tl{text}
        }   
     ,off .code:n =
        {
          \tl_set:Nn\l_tag_graphic_mode_tl{off}
        }   
     ,adjust-BBox .code:n = 
      {
        \bool_set_true:N \l_@@_graphic_bboxcorr_bool
        \seq_set_split:Nnn\l_@@_graphic_bboxcorr_seq{~}{#1~0pt~0pt~0pt~0pt}
      }
     ,
     tagging-setup .code:n= 
      {
        \keys_set:nn { tag/graphic }{#1}
      } 
%    \end{macrocode}
% only for legacy. Should perhaps be removed?
%    \begin{macrocode}
    ,tag .code:n = 
      {
        \str_case:nnF {#1}
          {
            {artifact}
            {
              \tl_set:Nn\l_tag_graphic_mode_tl{artifact}              
            }
            {false}{\tag_suspend:n{picture}}
          }
          {
            \tl_set:Nn\l_@@_graphic_struct_tl{#1}
            \tl_set:Nn\l_tag_graphic_mode_tl{alt}              
          }
      }     
   }          
%    \end{macrocode}
%
% \subsection{Tagging support for \cs{includegraphics}}
% 
% \subsubsection{User interface: Additional keys.}
%
% We also ensure that graphicx is loaded for the keyval support.
% At first a command to hold the tagging mode.
%    \begin{macrocode}
\tl_new:N  \l_tag_graphic_mode_tl
\tl_set:Nn \l_tag_graphic_mode_tl {alt} %TODO think about the right default.

\AddToHook{package/graphicx/after}[latex-lab]
  {
    \define@key{Gin}{alt}       
      {         
         \tl_if_empty:nF{#1}
           { \tl_set:Ne\l_@@_graphic_alt_tl{\text_purify:n{#1}}}
         \tl_set:Nn\l_tag_graphic_mode_tl {alt}
      }
    \define@key{Gin}{artifact}[]
      {
        \tl_set:Nn\l_tag_graphic_mode_tl{artifact}
      }
    \define@key{Gin}{actualtext}
     {
       \tl_if_empty:nF{#1}
         { \tl_set:Ne\l_@@_graphic_actual_tl{\text_purify:n{#1}} }   
       \tl_set:Nn\l_tag_graphic_mode_tl{actualtext}
     }
    \define@key{Gin}{tagging-setup}
     {
       \keys_set:nn { tag/graphic}{#1}
     }     
    \define@key{Gin}{adjust-BBox}
      {
        \bool_set_true:N \l_@@_graphic_bboxcorr_bool
        \seq_set_split:Nnn\l_@@_graphic_bboxcorr_seq{~}{#1~0pt~0pt~0pt~0pt}
      }
%    \end{macrocode}
% only for legacy, no longer documented
%    \begin{macrocode}
    \define@key{Gin}{tag}
      {
        \str_case:nnF {#1}
          {
            {artifact}
            {              
              \tl_set:Nn\l_tag_graphic_mode_tl{artifact}              
            }
            {false}{\tag_suspend:n{Gin}}
          }
          { 
            \tl_set:Nn\l_@@_graphic_struct_tl{#1}
            \tl_set:Nn\l_tag_graphic_mode_tl{alt}              
          }
      }  
  }
\AddToHook{package/graphics/after}[latex-lab]
  {\RequirePackage{graphicx}}  
%    \end{macrocode}
%
%
% \subsubsection{Patching graphics commands}
% All changes are currently done in \cs{Gin@setfile}.
% We mainly have to add the sockets in the right place and
% to rearrange a bit the \cs{ifGin@draft} test.
% 
%    \begin{macrocode}
\AddToHook{package/graphics/after}
 {
  \def\Gin@setfile#1#2#3{%  
  \ifx\\#2\\\Gread@false\fi
  \ifGin@bbox\else
    \ifGread@
      \csname Gread@%
         \expandafter\ifx\csname Gread@#1\endcsname\relax
           eps%
         \else
           #1%
         \fi
      \endcsname{\Gin@base#2}%
    \else
      \Gin@nosize{#3}%
    \fi
  \fi
  \Gin@viewport@code
  \Gin@nat@height\Gin@ury bp%
  \advance\Gin@nat@height-\Gin@lly bp%
  \Gin@nat@width\Gin@urx bp%
  \advance\Gin@nat@width-\Gin@llx bp%
  \Gin@req@sizes
  \expandafter\ifx\csname Ginclude@#1\endcsname\relax
    \Gin@drafttrue
    \expandafter\ifx\csname Gread@#1\endcsname\relax
      \@latex@error{Can not include graphics of type: #1}\@ehc
      \global\expandafter\let\csname Gread@#1\endcsname\@empty
    \fi
  \fi
  \leavevmode
%    \end{macrocode}
% Here the tagging begins. We want to catch also the draft box, 
% and for luatex tagging must be started before the \cs{setbox}.
%    \begin{macrocode}
  \tag_socket_use:ne {graphic/init} {tagging-setup=\l_tag_graphic_mode_tl}
  \tag_socket_use:nn {graphic/begin} {\Gin@base\Gin@ext} 
%    \end{macrocode}
% We store also the draft box in a box and do not output it directly
% so that we can calculate its BBox too. 
%    \begin{macrocode}
  \ifGin@draft
     \setbox\z@
      \hb@xt@\Gin@req@width{%
        \vrule\hss
        \vbox to \Gin@req@height{%
           \hrule \@width \Gin@req@width
           \vss
           \edef\@tempa{#3}%
           \rlap{ \ttfamily\expandafter\strip@prefix\meaning\@tempa}%
           \vss
           \hrule}%
        \hss\vrule}%
  \else
    \@addtofilelist{#3}%
    \ProvidesFile{#3}[Graphic~file~(type~#1)]%
    \setbox\z@\hbox{\csname Ginclude@#1\endcsname{#3}}%
    \dp\z@\z@
    \ht\z@\Gin@req@height
    \wd\z@\Gin@req@width
  \fi   
%    \end{macrocode}
% This ends the tagging. 
%    \begin{macrocode}
  \tag_socket_use:nnn{graphic/end}
    { \Gin@tag@bbox@attribute } 
    { \box\z@ } 
  }
 }  
%    \end{macrocode}
%
%
% \subsubsection{Calculating the BBox}
% 
% This is the large code part.
% 
% \begin{macro}{\@@_graphic_get_trim:}
% Graphics can be trimmed with the trim and the viewport key.
% If the graphic is not clipped the values must be taken into
% account when rotating.
% If viewport is used we have to calculate the trim.
% 
%    \begin{macrocode}
\cs_new_protected:Npn \@@_graphic_get_trim:
 { 
   \legacy_if:nTF {Gin@clip} 
%    \end{macrocode}
% Setting to 0 is not strictly needed but looks cleaner.
%    \begin{macrocode}
    {
      \fp_zero:N\l_@@_graphic_trim_lx_fp
      \fp_zero:N\l_@@_graphic_trim_ly_fp
      \fp_zero:N\l_@@_graphic_trim_ux_fp
      \fp_zero:N\l_@@_graphic_trim_uy_fp 
    } 
    { 
      \fp_set:Nn \l_@@_graphic_trim_lx_fp {\l_@@_graphic_scale_fp*\Gin@vllx}
      \fp_set:Nn \l_@@_graphic_trim_ly_fp {\l_@@_graphic_scale_fp*\Gin@vlly}
      \fp_set:Nn \l_@@_graphic_trim_ux_fp {\l_@@_graphic_scale_fp*\Gin@vurx}
      \fp_set:Nn \l_@@_graphic_trim_uy_fp {\l_@@_graphic_scale_fp*\Gin@vury}   
      \cs_if_exist:NT \Gin@ollx
        {
         \fp_set:Nn \l_@@_graphic_trim_ux_fp {\l_@@_graphic_scale_fp* (\Gin@ourx-(\Gin@urx)) }     
         \fp_set:Nn \l_@@_graphic_trim_uy_fp {\l_@@_graphic_scale_fp* (\Gin@oury-(\Gin@ury)) }     
        }
    }         
 }
%    \end{macrocode}
% \end{macro}
%
% \begin{macro}{\@@_graphic_get_scale:}
%    \begin{macrocode}
\cs_new_protected:Npn \@@_graphic_get_scale: 
 {
   \fp_set:Nn \l_@@_graphic_scale_fp
     {
       \str_if_eq:eeTF {\Gin@scalex} { ! }
        { \Gin@scaley }
        { \Gin@scalex }
     }
 }
%    \end{macrocode}
% \end{macro}
% 
% \begin{macro}{\@@_graphic_applyangle:nnnn}
% This takes the current BBox and rotates it according to the use angle.
% This is the most laborious code, as we have to take also the trim values into 
% account. We have to compare the values after the rotation to find the right corners
% for the BBox. Not sure, if this is the most effective code,
% the l3draw package has similar code to calculate a rotation,
% this can perhaps be reused ...
%    \begin{macrocode}
\cs_new_protected:Npn \@@_graphic_applyangle:nnnn #1#2#3#4 %lx,ly,ux,uy
 {
   \bool_lazy_and:nnT 
     {\cs_if_exist_p:N \Grot@angle } 
     {! \int_compare_p:nNn { \Grot@angle }={0}}
     { 
       \fp_set:Nn \l_@@_graphic_sin_fp  { sind(\Grot@angle) }
       \fp_set:Nn \l_@@_graphic_cos_fp  { cosd(\Grot@angle) }   
       \fp_set:Nn \l_@@_graphic_lx_fp {#1}    
       \fp_set:Nn \l_@@_graphic_ly_fp {#2}
       \fp_set:Nn \l_@@_graphic_ux_fp {#3}
       \fp_set:Nn \l_@@_graphic_uy_fp {#4}  
%    \end{macrocode}
%  get the x coordinates (cos,-sin)
%    \begin{macrocode}
       \fp_set:Nn\l_@@_graphic_lxly_fp 
        { 
          -\l_@@_graphic_trim_lx_fp *  \l_@@_graphic_cos_fp 
          +\l_@@_graphic_trim_ly_fp *  \l_@@_graphic_sin_fp 
        } 
       \fp_set:Nn\l_@@_graphic_lxuy_fp 
        { 
          (-\l_@@_graphic_trim_lx_fp) * \l_@@_graphic_cos_fp 
          +
           (\l_@@_graphic_uy_fp-\l_@@_graphic_ly_fp-\l_@@_graphic_trim_ly_fp) 
            * (-\l_@@_graphic_sin_fp) 
        } 
       \fp_set:Nn\l_@@_graphic_uxly_fp 
        { 
          (\l_@@_graphic_ux_fp-\l_@@_graphic_lx_fp-\l_@@_graphic_trim_lx_fp) 
           * \l_@@_graphic_cos_fp 
          +
          (\l_@@_graphic_trim_ly_fp) * (\l_@@_graphic_sin_fp) 
        }  
       \fp_set:Nn\l_@@_graphic_uxuy_fp 
        { 
          (\l_@@_graphic_ux_fp-\l_@@_graphic_lx_fp-\l_@@_graphic_trim_lx_fp) 
            *  \l_@@_graphic_cos_fp 
         +
          (\l_@@_graphic_uy_fp-\l_@@_graphic_ly_fp-\l_@@_graphic_trim_ly_fp) 
            * (-\l_@@_graphic_sin_fp) 
        }  
       \tl_gset:Ne\g_@@_graphic_lx_tl 
         { 
          \fp_eval:n
            {
              min 
               ( 
                 \l_@@_graphic_lxly_fp,
                 \l_@@_graphic_lxuy_fp, 
                 \l_@@_graphic_uxly_fp, 
                 \l_@@_graphic_uxuy_fp, 
               )
               +\l_@@_graphic_lx_fp
               +\l_@@_graphic_trim_lx_fp
            }
         }    
       \tl_gset:Ne\g_@@_graphic_ux_tl 
         { 
           \fp_eval:n
             {
               max 
                ( 
                  \l_@@_graphic_lxly_fp,
                  \l_@@_graphic_lxuy_fp, 
                  \l_@@_graphic_uxly_fp, 
                  \l_@@_graphic_uxuy_fp 
                )
                +\l_@@_graphic_lx_fp
                +\l_@@_graphic_trim_lx_fp                 
             }
         }
%    \end{macrocode}
% get the y coordinates (sin,cos)
%    \begin{macrocode}
       \fp_set:Nn\l_@@_graphic_lxly_fp 
        { 
           -\l_@@_graphic_trim_lx_fp *  \l_@@_graphic_sin_fp 
           -\l_@@_graphic_trim_ly_fp *  \l_@@_graphic_cos_fp  
        }   
       \fp_set:Nn\l_@@_graphic_lxuy_fp 
        { 
          - \l_@@_graphic_trim_lx_fp *  \l_@@_graphic_sin_fp 
          +
          (\l_@@_graphic_uy_fp-\l_@@_graphic_ly_fp-\l_@@_graphic_trim_ly_fp) 
            * \l_@@_graphic_cos_fp 
        } 
       \fp_set:Nn\l_@@_graphic_uxly_fp 
        { 
          (\l_@@_graphic_ux_fp-\l_@@_graphic_lx_fp-\l_@@_graphic_trim_lx_fp) 
            * \l_@@_graphic_sin_fp 
          - \l_@@_graphic_trim_ly_fp * \l_@@_graphic_cos_fp 
        }  
       \fp_set:Nn\l_@@_graphic_uxuy_fp 
        { 
          (\l_@@_graphic_ux_fp-\l_@@_graphic_lx_fp-\l_@@_graphic_trim_lx_fp) 
            * \l_@@_graphic_sin_fp 
         +
          (\l_@@_graphic_uy_fp-\l_@@_graphic_ly_fp-\l_@@_graphic_trim_ly_fp) 
            * \l_@@_graphic_cos_fp 
        }   
       \tl_gset:Ne\g_@@_graphic_ly_tl 
         { 
          \fp_eval:n
            {
              min 
               ( 
                 \l_@@_graphic_lxly_fp,
                 \l_@@_graphic_lxuy_fp, 
                 \l_@@_graphic_uxly_fp, 
                 \l_@@_graphic_uxuy_fp 
               ) 
              + \l_@@_graphic_ly_fp + \l_@@_graphic_trim_ly_fp 
            }
         }        
       \tl_gset:Ne\g_@@_graphic_uy_tl 
         { 
           \fp_eval:n
             {
               max 
                ( 
                  \l_@@_graphic_lxly_fp,
                  \l_@@_graphic_lxuy_fp, 
                  \l_@@_graphic_uxly_fp, 
                  \l_@@_graphic_uxuy_fp, 
                ) 
               + \l_@@_graphic_ly_fp + \l_@@_graphic_trim_ly_fp 
             }
         }           
     }      
 }  
\cs_generate_variant:Nn\@@_graphic_applyangle:nnnn {VVVV}
%    \end{macrocode}
% \end{macro}
% \begin{macro}{\@@_graphic_applycorr:NNNN}
% This command is used to add at the end the correction values. Quite dump ...
%    \begin{macrocode}
\cs_new_protected:Npn \@@_graphic_applycorr:NNNN #1 #2 #3 #4
 {
  \bool_if:NT\l_@@_graphic_bboxcorr_bool
    {
     \tl_set:Ne #1
       { 
         \fp_eval:n 
          {
           #1
           + 
           \dim_to_decimal_in_bp:n {\seq_item:Nn \l_@@_graphic_bboxcorr_seq {1} } 
          }
       }
     \tl_set:Ne #2
       { 
         \fp_eval:n 
          {
           #2
           + 
           \dim_to_decimal_in_bp:n {\seq_item:Nn \l_@@_graphic_bboxcorr_seq {2} } 
          }
       }
     \tl_set:Ne #3
       { 
         \fp_eval:n 
          {
           #3
           + 
           \dim_to_decimal_in_bp:n {\seq_item:Nn \l_@@_graphic_bboxcorr_seq {3} } 
          }
       }
     \tl_set:Ne #4
       { 
         \fp_eval:n 
          {
           #4
           + 
           \dim_to_decimal_in_bp:n {\seq_item:Nn \l_@@_graphic_bboxcorr_seq {4} } 
          }
       } 
     }  
  }              
%    \end{macrocode}
% \end{macro}
% 
% \begin{macro}[no-user-doc]{\Gin@tag@bbox@attribute}
% This is the main command to calculate and set the Bbox attribute of the 
% \cs{includegraphics} command. It also sets the reference point on the page.
% 
%    \begin{macrocode}
\cs_new_protected:Npn \Gin@tag@bbox@attribute
  {  
    \__tag_graphic_get_scale:
    \@@_graphic_get_trim:
    \int_gincr:N\g_@@_graphic_int
    \tl_set:Ne\l_@@_graphic_currentlabel_tl {_@@_graphic.\int_use:N \g_@@_graphic_int}
    \@@_graphic_savepos:e { \l_@@_graphic_currentlabel_tl }
    \tl_gset:Ne\g_@@_graphic_lx_tl
      {
        \dim_to_decimal_in_bp:n 
          { \property_ref:een {\l_@@_graphic_currentlabel_tl}{xpos}{0}sp }
      }  
    \tl_gset:Ne\g_@@_graphic_ly_tl
      {
        \dim_to_decimal_in_bp:n 
          { \property_ref:een {\l_@@_graphic_currentlabel_tl}{ypos}{0}sp }
      } 
    \tl_gset:Ne\g_@@_graphic_ux_tl
        {
          \fp_eval:n 
           { 
             \g_@@_graphic_lx_tl
             + 
             \dim_to_decimal_in_bp:n { \Gin@req@width }
           }  
        }         
     \tl_gset:Ne\g_@@_graphic_uy_tl
        {
          \fp_eval:n 
           {
             \g_@@_graphic_ly_tl
             + 
              \dim_to_decimal_in_bp:n  { \Gin@req@height }
           }  
        }
%    \end{macrocode}
% If the graphics is not clipped we must add the trim values.
%    \begin{macrocode}     
    \legacy_if:nF {Gin@clip}          
      {            
        \tl_gset:Ne\g_@@_graphic_ux_tl
          {
            \fp_eval:n 
              {
                \g_@@_graphic_ux_tl                       
                + 
                \l_@@_graphic_trim_ux_fp
              }  
          }
        \tl_gset:Ne\g_@@_graphic_lx_tl
          {
            \fp_eval:n 
              {
                \g_@@_graphic_lx_tl              
                 - 
                \l_@@_graphic_trim_lx_fp
              }  
          }  
        \tl_gset:Ne\g_@@_graphic_uy_tl
          {
            \fp_eval:n 
              {
                \g_@@_graphic_uy_tl           
                + 
                \l_@@_graphic_trim_uy_fp
              }  
          }                                   
        \tl_gset:Ne\g_@@_graphic_ly_tl
          {
            \fp_eval:n 
              {
                \g_@@_graphic_ly_tl              
                - 
                \l_@@_graphic_trim_ly_fp
              }  
          }                                          
      } 
%    \end{macrocode}
% If there is an angle we now rotate the values.
%    \begin{macrocode}
    \@@_graphic_applyangle:VVVV 
      \g_@@_graphic_lx_tl
      \g_@@_graphic_ly_tl
      \g_@@_graphic_ux_tl
      \g_@@_graphic_uy_tl 
%    \end{macrocode}
% At last we have to add the correction values
%    \begin{macrocode}
    \@@_graphic_applycorr:NNNN
      \g_@@_graphic_lx_tl
      \g_@@_graphic_ly_tl
      \g_@@_graphic_ux_tl
      \g_@@_graphic_uy_tl 
%    \end{macrocode}
%    \begin{macrocode}
  \bool_if:NT\l_@@_graphic_debug_bool
   {
     \@@_graphic_show_bbox:VVVVne
      \g_@@_graphic_lx_tl
      \g_@@_graphic_ly_tl
      \g_@@_graphic_ux_tl
      \g_@@_graphic_uy_tl
      {red}
      {\l_@@_graphic_currentlabel_tl}
   }
%    \end{macrocode}
% Now we add the attribute. We do it manually as it had to be delayed until now.
% The structure and the mc must be open earlier, before the \cs{setbox}  (at least
% for luatex it has to). 
%    \begin{macrocode}
    \tag_struct_gput:ene{\tag_get:n{struct_num}}{attribute}
       {
         /O /Layout /BBox~
           [
            \g_@@_graphic_lx_tl\c_space_tl
            \g_@@_graphic_ly_tl\c_space_tl
            \g_@@_graphic_ux_tl\c_space_tl
            \g_@@_graphic_uy_tl
           ]
       }          
  } 
%    \end{macrocode}
% \end{macro}
% 
% \subsection{Support for the picture environment}
% 
% \subsubsection{User interface}
% 
% The original picture has no key-val argument yet. In the new optional
% argument uses the generic tag/graphic keys.
% We could perhaps use the \texttt{Gin} keys instead, but there could
% be side-effects if some uses other \texttt{Gin} keys like 
% \texttt{angle}, so better stay on the safe side.
%
%
% \subsubsection{Calculation of the BBox}
% 
% \begin{macro}[no-user-doc]{\picture@tag@bbox@attribute}
% Picture needs a similar command to calculate the BBox. But here we stay simple
% and use simply the size of the picbox.
% 
%    \begin{macrocode}
\newcommand\picture@tag@bbox@attribute
 {
   \int_gincr:N\g_@@_graphic_int
   \tl_set:Ne\l_@@_graphic_currentlabel_tl {_@@_graphic.\int_use:N \g_@@_graphic_int}
   \@@_graphic_savepos:e { \l_@@_graphic_currentlabel_tl }
   \tl_gset:Ne \g_@@_graphic_lx_tl 
     { 
       \dim_to_decimal_in_bp:n 
       { \property_ref:een {\l_@@_graphic_currentlabel_tl}{xpos}{0}sp  } 
     }
   \tl_gset:Ne \g_@@_graphic_ly_tl 
     { 
       \dim_to_decimal_in_bp:n 
       { \property_ref:een {\l_@@_graphic_currentlabel_tl}{ypos}{0}sp - \dp\@picbox } 
     }  
   \tl_gset:Ne \g_@@_graphic_ux_tl 
     { 
       \dim_to_decimal_in_bp:n 
        {
          \g_@@_graphic_lx_tl bp + \wd\@picbox        
        }  
     }     
   \tl_gset:Ne \g_@@_graphic_uy_tl 
     { 
       \dim_to_decimal_in_bp:n 
        {
          \g_@@_graphic_ly_tl bp + \ht\@picbox + \dp\@picbox         
        }  
     }     
   \@@_graphic_applycorr:NNNN
        \g_@@_graphic_lx_tl
        \g_@@_graphic_ly_tl
        \g_@@_graphic_ux_tl
        \g_@@_graphic_uy_tl        
   \bool_if:NT\l_@@_graphic_debug_bool
     {
       \@@_graphic_show_bbox:VVVVne
        \g_@@_graphic_lx_tl
        \g_@@_graphic_ly_tl
        \g_@@_graphic_ux_tl
        \g_@@_graphic_uy_tl
        {red}
        {\l_@@_graphic_currentlabel_tl}
     }
%    \end{macrocode}
% this stores the attribute in the structure.
%    \begin{macrocode}
   \tag_struct_gput:ene{\tag_get:n{struct_num}}{attribute}
         {
             /O /Layout /BBox~
             [
              \g_@@_graphic_lx_tl\c_space_tl
              \g_@@_graphic_ly_tl\c_space_tl
              \g_@@_graphic_ux_tl\c_space_tl
              \g_@@_graphic_uy_tl 
             ]
        }            
  }
%    \end{macrocode}
% \end{macro}
%  
% \subsubsection{Patching the commands}
% We redefine \cs{picture} to accept an optional argument.
% We also ensure that we are in 
% hmode, so that stopping tagging doesn't confuse the paratags.
%    \begin{macrocode}
\RenewDocumentCommand\picture{O{}m}
 { 
   \leavevmode
   \tag_socket_use:nn{graphic/init}{#1}   
   \pictur@#2
 } 
%    \end{macrocode}
% inside the picture box we stop tagging. 
%    \begin{macrocode}
\def\@picture(#1,#2)(#3,#4){%
  \@defaultunitsset\@picht{#2}\unitlength
  \@defaultunitsset\@tempdimc{#1}\unitlength
  \tag_socket_use:nn{graphic/begin}{picture~environment}
  \tag_suspend:n{\@picture} %do not tag inside the picture box
  \setbox\@picbox\hb@xt@\@tempdimc\bgroup    
    \@defaultunitsset\@tempdimc{#3}\unitlength
    \hskip -\@tempdimc
    \@defaultunitsset\@tempdimc{#4}\unitlength
    \lower\@tempdimc\hbox\bgroup
      \ignorespaces}
%    \end{macrocode}
% 
%    \begin{macrocode}
\def\endpicture{%
  \egroup\hss\egroup
    \ht\@picbox\@picht\dp\@picbox\z@
    \tag_resume:n{\@picture}
    \tag_socket_use:nnn{graphic/end}
      {\picture@tag@bbox@attribute}
      {\mbox{\box\@picbox}}}
%    \end{macrocode}
%
%
% \subsection{Debugging code}
% This command puts a transparent layer in the size of the BBox over an graphic. 
% \begin{macro}{\@@_graphic_show_bbox:nnnnnn}
%    \begin{macrocode}
\cs_new_protected:Npn \@@_graphic_show_bbox:nnnnnn #1#2#3#4#5#6%#5 color, #6 label name
 {  
    \iow_log:n {tag/graphic~debug:~BBox~of~graphics~#6~is~#1~#2~#3~#4}  
    \hook_gput_code:nnn
     {shipout/foreground}
     {tag/graphic} 
     {
      \int_compare:nNnT 
       {\g_shipout_readonly_int}
       =
       {\property_ref:een{#6}{abspage}{0}}
       {
        \put
         (#1 bp,\dim_eval:n{-\paperheight + \dim_eval:n{#2 bp}})
         {
          \opacity_select:n{0.5}\color_select:n{#5}
          \rule
           {\dim_eval:n {#3 bp-\dim_eval:n{#1 bp}}}
           {\dim_eval:n {#4 bp-\dim_eval:n{#2 bp}}}
         }
       }           
     }
 }
 \cs_generate_variant:Nn  \@@_graphic_show_bbox:nnnnnn {VVVVne,nnnnne}
%    \end{macrocode}
% \end{macro}

%    \begin{macrocode}
%</package>  
%    \end{macrocode}

%    \begin{macrocode}
%<*latex-lab>
\ProvidesFile{graphic-latex-lab-testphase.ltx}
        [\ltlabgraphicdate\space v\ltlabgraphicversion\space latex-lab wrapper graphic]
\RequirePackage{latex-lab-testphase-graphic}
%</latex-lab>
%    \end{macrocode}
% \end{implementation}
